\chapter{The Four Fundamental Subspaces}

\section{Definitions}

\begin{definition}
Let $A$ be an $m\times n$ matrix. The four \emph{fundamental subspaces} of $A$ are
\begin{align*}
  \Col(A) &= \{A\vec x:\vec x\in\R^n\} = \Span\{\text{columns of }A\}\subseteq\R^m &&\text{(column space)}\\
  \Row(A) &= \Col(A^T) = \Span\{\text{rows of }A\}\subseteq\R^n &&\text{(row space)}\\
  \Null(A) &= \{\vec x\in\R^n : A\vec x=\vec 0\}\subseteq\R^n &&\text{(null space)}\\
  \Null(A^T) &= \{\vec y\in\R^m : A^T\vec y=\vec 0\}\subseteq\R^m &&\text{(left null space)}
\end{align*}
\end{definition}

If $L(\vec x)=A\vec x$ then $\Col(A)=\Range(L)$ and $\Null(A)=\Ker(L)$. All four are subspaces (they
are spans or solution spaces of homogeneous systems).

\section{Dimension of a subspace of $\R^n$}

\begin{theorem}
If $S$ is a non-trivial subspace of $\R^n$, then $S$ has a basis, and every basis of $S$ has the same
number of vectors.
\end{theorem}
\begin{proof}[Proof sketch]
Start with a non-zero vector of $S$ and keep adding vectors of $S$ not in the current span; each
addition keeps the set independent (Chapter 1), and by Chapter 3 an independent set in $\R^n$ has
at most $n$ vectors, so the process stops with a basis. If $\{\vec v_1,\dots,\vec v_k\}$ spans $S$
and $\{\vec w_1,\dots,\vec w_\ell\}\subseteq S$ with $\ell>k$, write each $\vec w_j$ in terms of the
$\vec v_i$: $W = VC$ with $C$ a $k\times\ell$ matrix. Since $\ell>k$ there is $\vec c\neq\vec0$
with $C\vec c=\vec0$, so $W\vec c=\vec0$ and the $\vec w_j$ are dependent. Applying this both ways shows
two bases have the same size.
\end{proof}

\begin{definition}[Dimension]
The \emph{dimension} $\dim S$ of a subspace $S$ is the number of vectors in any basis; $\dim\{\vec0\}=0$.
\end{definition}

\begin{theorem}
If $\dim S=k$, then any $k$ linearly independent vectors in $S$ form a basis, and any $k$ vectors that
span $S$ form a basis. If $S\subseteq T$ are subspaces then $\dim S\le\dim T$, with equality iff $S=T$.
\end{theorem}

\section{Finding bases}

The key observation: row operations change the columns of $A$, but they preserve every linear
relation among the columns (since $A\vec c=\vec0$ iff $R\vec c=\vec0$ when $A$ and $R$ are row
equivalent). They also do not change the row space (new rows are combinations of old rows and vice
versa).

\begin{theorem}
Let $R$ be the RREF of $A$ (or any REF).
\begin{enumerate}
  \item The pivot columns \emph{of $A$} (not of $R$) form a basis for $\Col(A)$.
  \item The non-zero rows of $R$ form a basis for $\Row(A)$.
  \item The vectors obtained from the parametric vector form of the solution to $A\vec x=\vec0$
        (one per free variable) form a basis for $\Null(A)$.
\end{enumerate}
Consequently $\dim\Col(A)=\dim\Row(A)=\rank A$ and $\dim\Null(A)=n-\rank A$.
\end{theorem}
\begin{proof}[Proof sketch]
(1) In $R$ the pivot columns are distinct standard basis vectors, hence independent, and every
non-pivot column is a combination of them. These relations transfer to the columns of $A$.
(2) The non-zero rows of $R$ span $\Row(R)=\Row(A)$ and are independent because each has a leading one
in a column where the others have zeros. (3) Each basic vector has a $1$ in its own free variable's
position and $0$ in the other free positions, giving independence; spanning is the parametric
form itself.
\end{proof}

\begin{important}[Pitfall]
Row operations \emph{change} the column space. The columns of $R$ tell you \emph{which} columns
to take, but the basis vectors must come from $A$ itself. Conversely, for the row space you may use
the rows of $R$ (and the rows of $A$ in the pivot positions are \emph{not} necessarily a basis if
you swapped rows).
\end{important}

\begin{example}
Let $A=\begin{bmatrix}1&2&0&1\\2&4&1&4\\1&2&1&3\end{bmatrix}$. As in Chapter 3,
\[
A\xrightarrow{\substack{R_2-2R_1\\R_3-R_1}}
\begin{bmatrix}1&2&0&1\\0&0&1&2\\0&0&1&2\end{bmatrix}
\xrightarrow{R_3-R_2}
\begin{bmatrix}1&2&0&1\\0&0&1&2\\0&0&0&0\end{bmatrix}=R .
\]
$\rank A=2$, with pivots in columns $1$ and $3$.
\begin{itemize}
  \item $\Col(A)$: basis $\{(1,2,1),(0,1,1)\}$ (columns 1 and 3 of $A$), dimension $2$.
  \item $\Row(A)$: basis $\{(1,2,0,1),(0,0,1,2)\}$, dimension $2$.
  \item $\Null(A)$: $x_2=s$, $x_4=t$, $x_1=-2s-t$, $x_3=-2t$, so the basis is
        $\{(-2,1,0,0),(-1,0,-2,1)\}$, dimension $2=4-2$.
  \item $\Null(A^T)$: row reduce $A^T$:
  \[
  \begin{bmatrix}1&2&1\\2&4&2\\0&1&1\\1&4&3\end{bmatrix}
  \xrightarrow{\substack{R_2-2R_1\\R_4-R_1}}
  \begin{bmatrix}1&2&1\\0&0&0\\0&1&1\\0&2&2\end{bmatrix}
  \xrightarrow{\substack{R_4-2R_3\\R_2\leftrightarrow R_3}}
  \begin{bmatrix}1&2&1\\0&1&1\\0&0&0\\0&0&0\end{bmatrix}
  \xrightarrow{R_1-2R_2}
  \begin{bmatrix}1&0&-1\\0&1&1\\0&0&0\\0&0&0\end{bmatrix}.
  \]
  So $y_3=t$, $y_1=t$, $y_2=-t$: basis $\{(1,-1,1)\}$, dimension $1=3-2$.
\end{itemize}
Sanity check: $(1,-1,1)$ is orthogonal to both basis vectors of $\Col(A)$: $1-2+1=0$ and
$0-1+1=0$. Likewise $(1,2,0,1)\cdot(-2,1,0,0)=0$, $(1,2,0,1)\cdot(-1,0,-2,1)=0$,
$(0,0,1,2)\cdot(-2,1,0,0)=0$, $(0,0,1,2)\cdot(-1,0,-2,1)=0$.
\end{example}

\section{The Rank--Nullity Theorem}

\begin{definition}
The \emph{nullity} of $A$ is $\nullity A=\dim\Null(A)$.
\end{definition}

\begin{theorem}[Rank--Nullity / Dimension Theorem]
For any $m\times n$ matrix $A$,
\[
  \rank A + \nullity A = n .
\]
Equivalently, for linear $L:\R^n\to\R^m$, $\dim\Range(L)+\dim\Ker(L)=n$.
\end{theorem}
\begin{proof}
In the RREF, each of the $n$ columns is either a pivot column (there are $\rank A$ of them) or
corresponds to a free variable (each contributing one basis vector of $\Null(A)$).
\end{proof}

\begin{theorem}
$\rank A = \rank A^T$. Hence $\dim\Null(A^T)=m-\rank A$.
\end{theorem}
\begin{proof}
$\rank A^T=\dim\Col(A^T)=\dim\Row(A)=\rank A$.
\end{proof}

\begin{example}
A $5\times 7$ matrix $A$ with $\rank A=3$ has $\dim\Col A=\dim\Row A=3$, $\nullity A=7-3=4$ and
$\dim\Null(A^T)=5-3=2$. It cannot have rank $6$, since $\rank A\le\min(5,7)$.
\end{example}

\section{Orthogonality of the fundamental subspaces}

\begin{definition}
For a subspace $S$ of $\R^n$, its \emph{orthogonal complement} is
$S^\perp=\{\vec x\in\R^n:\vec x\cdot\vec s=0 \text{ for all } \vec s\in S\}$.
\end{definition}

\begin{theorem}
$\Null(A)=\Row(A)^\perp$ in $\R^n$, and $\Null(A^T)=\Col(A)^\perp$ in $\R^m$.
\end{theorem}
\begin{proof}
$A\vec x=\vec0$ says exactly that $\vec x$ has dot product zero with every row of $A$, hence with
every linear combination of rows. Apply this to $A^T$ for the second statement.
\end{proof}

Together with rank--nullity this says $\R^n$ is split into two orthogonal pieces of dimensions
$\rank A$ and $n-\rank A$, and similarly $\R^m$. This also gives a test for consistency:
$A\vec x=\vec b$ is consistent iff $\vec b\in\Col(A)$ iff $\vec b$ is orthogonal to every vector in
$\Null(A^T)$. In the example above, $A\vec x=\vec b$ is consistent iff $b_1-b_2+b_3=0$.

\section*{Summary}
\begin{itemize}
  \item $\Col A$, $\Row A$, $\Null A$, $\Null A^T$; $\Col A=\Range$, $\Null A=\Ker$.
  \item Basis of $\Col A$: pivot columns of $A$. Basis of $\Row A$: non-zero rows of the RREF.
        Basis of $\Null A$: vectors from the free variables.
  \item $\dim\Col A=\dim\Row A=\rank A$; rank--nullity: $\rank A+\nullity A=n$.
  \item $\Null A\perp\Row A$ and $\Null A^T\perp\Col A$.
\end{itemize}
